% Desarrollo basado en los apuntes del estudiante; no es una transcripción.
\par\textbf{\textcolor{unabblue}{\faLightbulb}\quad Más allá de la corrección del código}\par
Un sistema puede procesar datos correctamente en el backend y fracasar si el
usuario no comprende cómo utilizarlo. La corrección de la lógica interna no
garantiza que una persona encuentre una función, interprete un mensaje o complete
una tarea. Cuando la interfaz exige adivinar qué hacer, la capacidad técnica
no se convierte en una solución efectiva.

Desde la Ingeniería en Sistemas, verificar el resultado de una operación debe
complementarse con comprobar que el usuario pueda iniciarla y reconocer su
finalización. Una navegación confusa puede generar errores de uso y abandono,
aunque el código cumpla sus especificaciones. Comprender a las personas es,
por tanto, parte de la definición del problema y de la validación del producto.

\par\textbf{\textcolor{unabblue}{\faCogs}\quad Pilares y convergencia interdisciplinaria}\par
Los pilares son \textbf{tecnología, creatividad, diseño y experiencia}:
implementar funciones, explorar alternativas, organizar las interacciones y
valorar su utilidad. Su integración requiere colaboración entre diseño gráfico
y multimedia, computación e Ingeniería en Sistemas, con enfoque en UX.
La arquitectura, los flujos y las decisiones visuales deben responder al mismo
objetivo; el diseño no puede reducirse a decorar un producto ya construido.

\par\textbf{\textcolor{unabblue}{\faUsers}\quad UI y UX: funciones complementarias}\par
La \textbf{interfaz de usuario (UI)} comprende la capa visual y tangible:
iconos, gráficos, colores, tipografías y controles. La \textbf{experiencia de
usuario (UX)} abarca el recorrido cognitivo y funcional: comprender, encontrar,
realizar una acción, obtener ayuda y alcanzar un objetivo.

\clearpage
\subsubsection*{\textcolor{unabblue}{\faIcon{sync-alt}}\quad Ciclo de trabajo iterativo}
El recorrido de ocho pasos relaciona la exploración del problema con la
construcción y la evaluación de la solución. La disposición en S representa
su secuencia; el retorno del paso 7 al 2 expresa que la mejora requiere volver
a comprender al usuario, no limitarse a corregir la interfaz.

\begin{center}
\begin{tikzpicture}[x=4.7cm,y=1.8cm,
  etapa/.style={draw=unabblue,fill=softgray,rounded corners=3pt,
    text width=3.6cm,minimum height=1.05cm,align=center,
    font=\small\sffamily,inner sep=5pt},
  flujo/.style={->,thick,draw=unabblue,>=stealth}]
  \node[etapa] (a) at (0,0) {\textbf{1.} Identificar\\la necesidad};
  \node[etapa] (b) at (1,0) {\textbf{2.} Conocer\\al usuario};
  \node[etapa] (c) at (2,0) {\textbf{3.} Diseñar\\la propuesta};
  \node[etapa] (d) at (2,-1) {\textbf{4.} Crear\\el prototipo};
  \node[etapa] (e) at (1,-1) {\textbf{5.} Desarrollar};
  \node[etapa] (f) at (0,-1) {\textbf{6.} Probar\\con usuarios};
  \node[etapa] (g) at (0,-2) {\textbf{7.} Mejorar};
  \node[etapa] (h) at (1.5,-2) {\textbf{8.} Solución digital\\centrada en las personas};
  \draw[flujo] (a)--(b); \draw[flujo] (b)--(c);
  \draw[flujo] (c)--(d); \draw[flujo] (d)--(e);
  \draw[flujo] (e)--(f); \draw[flujo] (f)--(g);
  \draw[flujo] (g)--(h);
  \draw[flujo,dashed] (g.west)--(-0.62,-2)--(-0.62,0.65)
    --(1,0.65)--(b.north);
\end{tikzpicture}

{\small\textbf{Ciclo de creación.} La flecha discontinua indica el retorno
7 $\rightarrow$ 2. Elaboración propia a partir de los apuntes.\par}
\end{center}

\textbf{Comprender y proponer (pasos 1 a 3).} Identificar la necesidad delimita
el problema que se busca resolver. Conocer al usuario permite examinar sus
objetivos, dificultades y contexto de uso. Diseñar la propuesta transforma
esa comprensión en una organización de funciones e interacciones, evitando
que la solución dependa únicamente de supuestos del equipo técnico.

\textbf{Construir y contrastar (pasos 4 a 6).} El prototipo hace visible la
propuesta antes de desarrollar todos sus componentes. El desarrollo materializa
las decisiones de diseño, mientras que las pruebas con usuarios permiten
observar si las personas entienden el recorrido y completan las tareas.
Estas pruebas complementan las verificaciones del funcionamiento interno.

\textbf{Aprender y mejorar (pasos 7 y 8).} La mejora convierte lo observado
en ajustes y nuevas preguntas. Retornar al paso 2 permite revisar necesidades
mal interpretadas y reconsiderar la propuesta, el prototipo o el desarrollo.
El paso 8 representa una solución contrastada con las personas; no significa
que deje de necesitar evaluación cuando cambien sus necesidades.

\clearpage
\subsubsection*{\textcolor{unabblue}{\faCheckCircle}\quad El usuario como criterio de calidad}
Los atributos del usuario deben traducirse en requisitos de calidad observables
y verificables. La tabla propone criterios de revisión derivados de los apuntes,
no resultados medidos durante el congreso.

\begingroup
\small
\renewcommand{\arraystretch}{1.2}
\begin{tabularx}{\linewidth}{@{}p{3.2cm}X@{}}
\toprule
\textbf{Atributo} & \textbf{Aplicación como requisito de calidad} \\
\midrule
Claridad & Las etiquetas y los mensajes permiten comprender qué hace una
función y qué se espera del usuario, sin depender de términos técnicos internos. \\
Facilidad de uso & Las tareas siguen una secuencia comprensible. Las pruebas
permiten identificar pasos innecesarios y dificultades para completar acciones. \\
Accesibilidad & La interacción contempla distintas capacidades y formas de
acceso, por ejemplo, navegación con teclado, texto legible y señales que no
dependan exclusivamente del color. \\
Coherencia visual & Los controles, términos y patrones mantienen un significado
consistente para que lo aprendido en una pantalla sea útil en las demás. \\
Retroalimentación & El sistema comunica el progreso, el resultado de las
acciones y los errores, e indica cómo continuar o recuperar la tarea. \\
\bottomrule
\end{tabularx}
\endgroup

\subsubsection*{\textcolor{unabblue}{\faLaptopCode}\quad Aplicación en Ingeniería en Sistemas}
Estos criterios pueden incorporarse al levantamiento de requisitos, al prototipado
y a la validación de proyectos. Como ejemplo de aplicación académica, una
herramienta de administración de redes debe permitir comprender el alcance
de una operación y confirmar su resultado. Evaluar etiquetas, permisos visibles
y mensajes vincula el funcionamiento técnico con el uso responsable del sistema.

\subsubsection*{\textcolor{unabblue}{\faRobot}\quad Inteligencia artificial como apoyo transversal}
La inteligencia artificial actúa como catalizador creativo y técnico: puede
apoyar la exploración de alternativas, el prototipado y el desarrollo. Sus
propuestas requieren revisión según el contexto del usuario. Como apoyo
transversal al trabajo interdisciplinario y a la iteración, no sustituye la
validación con personas de una solución comprensible, accesible y útil.
